\chapter{ATIVIDADES DESENVOLVIDAS}
\label{cap:atividades}

Durante o período de estágio supervisionado, foi possível acompanhar
atendimentos e procedimentos clínicos em diversas áreas da medicina
equina, proporcionando uma experiência abrangente e enriquecedora para a
formação profissional.

A rotina diária tinha início com a administração das medicações
previstas para o horário, seguida da realização dos exames clínicos dos
animais internados, com aferição de parâmetros vitais fundamentais para
o monitoramento clínico e a detecção precoce de alterações. Os dados
obtidos eram registrados em fichas clínicas individuais,
disponibilizadas sobre a bancada principal do setor. Posteriormente,
eram executadas as demais atividades, que incluíam o manejo de feridas,
além do auxílio em procedimentos clínicos de maior complexidade conforme
a demanda dos pacientes internados.

No âmbito dos atendimentos externos, foi possível acompanhar a equipe
docente da instituição em três ocasiões distintas. A primeira delas
envolveu um deslocamento ao município de Abaeté, Minas Gerais, onde
foram realizados procedimentos odontológicos, avaliação do sistema
locomotor, aplicação de protocolos vacinais, entre outras práticas
clínicas, sob condução da docente responsável. As outras duas ocorreram
em Cachoeira do Campo, distrito de Ouro Preto e na própria cidade de
Belo Horizonte, sendo ambas voltadas à avaliação de membros de cavalos
de esporte por meio de ultrassonografia. Essas experiências foram
especialmente relevantes por proporcionarem contato com a rotina clínica
fora do ambiente hospitalar, ampliando a visão sobre as diferentes
realidades e demandas que envolvem a atuação do médico veterinário de
equinos a campo.

Ainda durante o estágio supervisionado, como forma de integração e
estímulo ao aprendizado coletivo, foi proposto pelos residentes do setor
que os estagiários conduzissem uma apresentação oral para os demais
integrantes do hospital que tivessem interesse em assistir. A acadêmica
abordou o tema acidente ofídico em equinos por serpentes do gênero
\textit{Bothrops}, assunto de relevância clínica na medicina equina, cuja
elaboração contribuiu para o aprofundamento do conhecimento acerca desse
atendimento emergencial.

Por fim, foi possível participar de treinamentos promovidos pelo Núcleo
de Estudos em Medicina Interna Equina (EMIE), que se tornaram momentos
muito enriquecedores de troca e aprendizado coletivo. Além disso, houve
a oportunidade de acompanhar aulas práticas das disciplinas de
Semiologia Veterinária e Clínica Médica de Grandes Animais, realizadas
no próprio setor, enriquecendo ainda mais a experiência do estágio
supervisionado.
